% ---------------------------------------------------------------
%  Manual del usuario de REM · Ramon Guillamon · CC BY-SA 4.0
% ---------------------------------------------------------------
\usepackage[a4paper,inner=2.8cm,outer=2.2cm,top=2.6cm,bottom=2.6cm]{geometry}
\usepackage{fontspec}
\usepackage[spanish,es-noshorthands]{babel}
\setmainfont{Palatino}
\setsansfont{Helvetica Neue}
\setmonofont[Scale=0.86]{Menlo}
\newfontfamily\symfont{Apple Symbols}
\newcommand{\cmd}{{\symfont ⌘}}
\usepackage{xcolor}
\definecolor{accent}{HTML}{0A6CD9}
\definecolor{accent2}{HTML}{0E7490}
\definecolor{mk}{HTML}{C2410C}
\definecolor{codebg}{HTML}{F4F6FA}
\definecolor{ideabg}{HTML}{EEF5FF}
\definecolor{warnbg}{HTML}{FEF2F2}
\definecolor{warnframe}{HTML}{DC2626}
\definecolor{plainbg}{HTML}{ECFEFF}
\definecolor{checkbg}{HTML}{ECFDF5}
\definecolor{checkframe}{HTML}{059669}
\definecolor{objbg}{HTML}{F5F3FF}
\definecolor{objframe}{HTML}{7C3AED}
\usepackage{microtype}
\usepackage{titlesec}
\usepackage{fancyhdr}
\usepackage{graphicx}
\usepackage{float}
\usepackage{booktabs}
\usepackage{longtable}
\usepackage{array}
\usepackage{calc}
\usepackage{enumitem}
\usepackage{tikz}
\usepackage[most]{tcolorbox}
\usepackage{fvextra}
\usepackage{caption}
\usepackage{hyperref}
\usepackage{bookmark}
\hypersetup{colorlinks=true,linkcolor=accent,urlcolor=accent,citecolor=accent,pdftitle={Manual del usuario de REM: gestión local para autónomos y pymes, sin cloud},pdfauthor={Ramon Guillamon},pdfsubject={Manual de usuario y curso de construcción de aplicaciones con IA},pdfkeywords={REM, autónomos, pymes, IA local, aplicación web, curso}}
\captionsetup{font={small,sf},labelfont={bf,color=accent},labelsep=period,justification=raggedright,singlelinecheck=false}
\setlength{\parindent}{0pt}\setlength{\parskip}{0.55em}
\linespread{1.08}
\titleformat{\part}[display]{\sffamily\bfseries\color{accent}\Huge}{Parte \thepart}{10pt}{\Huge}
\titleformat{\chapter}[display]{\sffamily\bfseries\color{accent}\huge}{\chaptertitlename\ \thechapter}{8pt}{\Huge}[\vspace{2pt}{\color{accent!40}\titlerule[1.2pt]}]
\titleformat{\section}{\sffamily\Large\bfseries\color{accent}}{\thesection}{0.7em}{}
\titleformat{\subsection}{\sffamily\large\bfseries\color{accent}}{\thesubsection}{0.7em}{}
\titleformat{\subsubsection}{\sffamily\normalsize\bfseries\color{accent2}}{\thesubsubsection}{0.7em}{}
\pagestyle{fancy}\fancyhf{}
\fancyhead[LE]{\sffamily\small\nouppercase{\leftmark}}
\fancyhead[RO]{\sffamily\small Manual del usuario de REM}
\fancyfoot[LE,RO]{\sffamily\small\thepage}
\renewcommand{\headrulewidth}{0.3pt}\setlength{\headheight}{14pt}
\fancypagestyle{plain}{\fancyhf{}\fancyfoot[C]{\sffamily\small\thepage}\renewcommand{\headrulewidth}{0pt}}
\setlist{itemsep=0.2em,topsep=0.35em}
\setcounter{tocdepth}{2}
% --- Recuadros ---
\newtcolorbox{ideabox}{colback=ideabg,colframe=accent,title=\textbf{Idea clave},arc=1mm,boxrule=0.6pt,left=2mm,right=2mm,top=1mm,bottom=1mm,fonttitle=\sffamily}
\newtcolorbox{warnbox}{colback=warnbg,colframe=warnframe,title=\textbf{Cuidado},arc=1mm,boxrule=0.6pt,left=2mm,right=2mm,top=1mm,bottom=1mm,fonttitle=\sffamily}
\newtcolorbox{plainbox}[1]{colback=plainbg,colframe=accent2,title=\textbf{En cristiano: #1},arc=1mm,boxrule=0.6pt,left=2mm,right=2mm,top=1mm,bottom=1mm,fonttitle=\sffamily}
\newtcolorbox{checkbox}{colback=checkbg,colframe=checkframe,title=\textbf{Pruébalo},arc=1mm,boxrule=0.6pt,left=2mm,right=2mm,top=1mm,bottom=1mm,fonttitle=\sffamily}
\newtcolorbox{objbox}{colback=objbg,colframe=objframe,title=\textbf{Objetivos de aprendizaje},arc=1mm,boxrule=0.6pt,left=2mm,right=2mm,top=1mm,bottom=1mm,fonttitle=\sffamily}
\newtcolorbox{aibox}{colback=ideabg,colframe=accent2,title=\textbf{Con la IA},arc=1mm,boxrule=0.6pt,left=2mm,right=2mm,top=1mm,bottom=1mm,fonttitle=\sffamily}
\DefineVerbatimEnvironment{codebox}{Verbatim}{frame=single,framerule=0.5pt,rulecolor=\color{accent!40},framesep=6pt,breaklines=true,breakanywhere=true,fontsize=\small,xleftmargin=1mm,xrightmargin=1mm,commandchars=\\\{\}}
% --- Marcas numeradas y capturas ---
\newcommand{\mk}[1]{\tikz[baseline=(c.base)]\node[circle,fill=mk,text=white,inner sep=0.8pt,minimum size=4.3mm,font=\sffamily\bfseries\scriptsize] (c) {#1};}
\newenvironment{leyenda}{\begin{description}[font=\normalfont,leftmargin=1.2cm,labelwidth=1cm,labelsep=0.1cm,align=left,itemsep=0.18em,topsep=0.3em]}{\end{description}}
\setlength{\fboxsep}{0pt}\setlength{\fboxrule}{0.5pt}
\newcommand{\figura}[3][0.60]{\begin{figure}[!htbp]\centering\fcolorbox{black!18}{white}{\includegraphics[width=\linewidth,height=#1\textheight,keepaspectratio]{figures/#2}}\caption{#3}\label{fig:#2}\end{figure}}
\newcommand{\movil}[4][0.52]{\begin{figure}[!htbp]\centering\foreach \f in {#2}{\fcolorbox{black!18}{white}{\includegraphics[height=#1\textheight]{figures/\f}}\hspace{0.012\linewidth}}\caption{#3}\label{fig:#4}\end{figure}}
\newcommand{\ruta}[1]{\texttt{\small #1}}
\newcommand{\ui}[1]{\textsf{\textbf{#1}}}
% --- compatibilidad con la salida de pandoc ---
\providecommand{\tightlist}{\setlength{\itemsep}{0pt}\setlength{\parskip}{0pt}}
\providecommand{\real}[1]{#1}
\providecommand{\pandocbounded}[1]{#1}
\providecommand{\LTcaptype}{none}
\newcounter{none}
\RecustomVerbatimEnvironment{verbatim}{Verbatim}{frame=single,framerule=0.4pt,rulecolor=\color{accent!35},framesep=5pt,breaklines=true,breakanywhere=true,fontsize=\footnotesize,xleftmargin=1mm,xrightmargin=1mm}
\usepackage{newunicodechar}
\newunicodechar{→}{{\symfont →}}
\newunicodechar{←}{{\symfont ←}}
\newunicodechar{≈}{{\symfont ≈}}
\newunicodechar{✓}{{\symfont ✓}}
\newunicodechar{✔}{{\symfont ✔}}
\newunicodechar{⌘}{{\symfont ⌘}}
\newunicodechar{▸}{{\symfont ▸}}
\renewcommand{\topfraction}{0.9}\renewcommand{\bottomfraction}{0.9}\renewcommand{\textfraction}{0.08}\renewcommand{\floatpagefraction}{0.7}\setcounter{topnumber}{3}\setcounter{totalnumber}{4}
